\section{Contributions}

This thesis has delivered a set of interrelated contributions that advance the state of the art in autonomous marine robotics. The work integrates simulation development, planning and control strategies, and data-driven modeling approaches, all directed at improving the reliability, safety, and adaptability of \ac{MPC}-based systems for underwater inspection missions.

\noindent \textbf{Visually Realistic Simulation Platform}  
A key foundational contribution is the development of \textit{UNav-Sim}, a high-fidelity, open-source underwater robotics simulator leveraging Unreal Engine 5. Unlike conventional simulation tools that often sacrifice either visual realism or dynamic fidelity, \textit{UNav-Sim} uniquely combines photorealistic rendering of underwater environments with accurate hydrodynamic modeling of robot motion. This capability substantially narrows the simulation-to-reality gap, providing a robust platform for the development, benchmarking, and validation of perception and control algorithms before deployment in real-world operations.

\vspace{0.75em}

\noindent \textbf{Entanglement-Aware Coverage Planning}  
The research introduces the \textit{REACT} framework, a novel real-time coverage planning approach explicitly designed to address tether management for remotely operated vehicles in cluttered inspection environments. This framework incorporates a geometry-aware tether model into the planning loop, enabling proactive avoidance of entanglement and automatic reactive replanning in dynamic scenarios. By systematically balancing coverage completeness with tether safety constraints, this contribution enables high-assurance inspection missions in environments that were previously impractical due to entanglement risks.

\vspace{0.75em}

\noindent \textbf{Perception-Aware Nonlinear MPC}  
This thesis advances the capabilities of model predictive control by embedding task-specific perception objectives into the control optimization process. Through the \textit{VT-NMPC} approach, the system dynamically selects robot configurations that maintain optimal viewpoints for visual inspection while respecting physical constraints and environmental disturbances. This integration departs from traditional trajectory tracking methods by explicitly optimizing inspection quality alongside navigation, resulting in higher mission effectiveness and robustness.

\vspace{0.75em}

\noindent \textbf{Physics-Informed Neural Modeling of Dynamics}  
To improve predictive accuracy critical for MPC, the research presents the application of \textit{PINC} for learning full dynamic models of underwater vehicles from experimental and simulation data. These models respect governing physical laws, resulting in enhanced generalizability across operating conditions compared to purely data-driven or purely analytical models. This contribution demonstrates how hybrid approaches (combining data-driven methods with first-principles modeling) can reduce modeling error and improve long-horizon prediction fidelity, which is essential for safe and efficient control.

\vspace{0.75em}

\noindent \textbf{Adaptive Online Learning of Environmental Disturbances}  
Recognizing the significant impact of unmodeled and time-varying disturbances on control performance, this work develops \textit{DF-GP}, an adaptive Gaussian Process framework for real-time disturbance estimation. By continuously learning residual dynamics and disturbance forces during operation, this method allows the controller to compensate adaptively without requiring manual retuning. The result is a substantial improvement in trajectory tracking accuracy and operational robustness across diverse and changing environmental conditions.

\vspace{0.75em}

\noindent Collectively, these contributions represent a comprehensive step toward more intelligent, reliable, and adaptable marine robotic systems. By systematically addressing challenges spanning simulation, planning, control, and learning, this thesis demonstrates how integrated frameworks can elevate the autonomy and mission success of underwater inspection robots operating in complex and uncertain environments.
